% -*- TeX-master: "../main.tex" -*-
\chapter{Wybór modelu autobusu}
Zaprojektowanie sieci przystanków i rozkładów jazdy musi iść w parze z odpowiednim doborem taboru, który fizycznie obsłuży zaplanowane trasy. Wybór modelu autobusu jest złożonym problemem decyzyjnym, wymagającym uwzględnienia ograniczeń budżetowych, wymagań terenowych oraz oczekiwań pasażerów. Ponieważ żaden dostępny na rynku pojazd nie jest idealny pod każdym względem (np. pojazdy najtańsze nie zawsze oferują najwyższy prestiż lub odpowiednią moc), konieczne było zastosowanie narzędzi z zakresu badań operacyjnych.

W niniejszym rozdziale przeprowadzono ocenę dostępnych na rynku autobusów z wykorzystaniem metody SAW (\textit{Simple Additive Weighting}). Pozwala ona na sprowadzenie różnorodnych parametrów technicznych i ekonomicznych do wspólnego mianownika, ułatwiając tym samym obiektywny wybór najlepszego rozwiązania dla Sławkowa.

\section{Założenia metody i wagi kryteriów}
Wychodząc naprzeciw specyficznym potrzebom nowo projektowanego układu komunikacyjnego w Sławkowie, parametry modelu decyzyjnego dostosowano w sposób odzwierciedlający realne wyzwania topograficzne oraz politykę transportową miasta. W procesie zrezygnowano z równego podziału wag, silnie premiując cechy gwarantujące sprawność obsługi i zadowolenie podróżnych. Wyodrębniono następujące kryteria:

\begin{itemize}
    \item \textbf{K1: Moc silnika [KM] (stymulanta, waga $0,26$)} -- podwyższona waga tego kryterium wynika bezpośrednio z topografii miasta Sławków. Nierówny teren i przewyższenia sprawiają, że odpowiednia moc napędowa jest absolutnie kluczowa do sprawnego pokonywania wzniesień i utrzymania płynności rozkładu jazdy.
    \item \textbf{K2: Miejsca dla pasażerów (stymulanta, waga $0,26$)} -- wysoka waga ma na celu zminimalizowanie zjawiska przepełnienia i "ciśnienia się" pasażerów w godzinach szczytowych, co bezpośrednio przekłada się na chęć korzystania z transportu publicznego.
    \item \textbf{K3: Prestiż i komfort [skala 1-10] (stymulanta, waga $0,26$)} -- subiektywna ocena uśredniona przez grono decydentów w sekcji \ref{sec:bus_prestige}. Uwzględnia wygląd zewnętrzny, estetykę wnętrza i ciekawe rozwiązania technologiczne. Priorytetyzacja tego aspektu wynika z chęci zapewnienia mieszkańcom jak najlepszych doświadczeń z podróży.
    \item \textbf{K4: Udogodnienia dla niepełnosprawnych (stymulanta, waga $0,11$)} -- kryterium dyskretne zliczające łączną ilość zastosowanych w pojeździe ułatwień dla osób o ograniczonej mobilności (np. rampa, niska podłoga, przyklęk, miejsce na wózek, braille).
    \item \textbf{K5: Koszt zakupu [mln PLN] (destymulanta, waga $0,11$)} -- szacunkowy koszt jednostkowy pojazdu.
\end{itemize}

Aby zapewnić miarodajność, analizę rozdzielono na dwie niezależne kategorie rynkowe: autobusy duże (przegubowe i wielkopojemne) oraz autobusy średnie (pojazdy klasy MIDI).

\section{Ocena prestiżu autobusów}\label{sec:bus_prestige}

Jednym z kluczowych kryteriów przyjętych w analizie decyzyjnej jest prestiż pojazdu. Wartość tego parametru wyznaczono jako średnią arytmetyczną z ocen eksperckich zespołu projektowego. Składowe oceny obejmowały analizę designu zewnętrznego, aranżacji wnętrza, a także identyfikację niestandardowych udogodnień i nowatorskich rozwiązań konstrukcyjnych.

Dodatkowym atutem, który bezpośrednio wpłynął na podwyższenie punktacji dla pojazdów marki Solaris, jest obecność charakterystycznej maskotki jamnika. Zespół decyzyjny docenił ten nietypowy element wizerunkowy, ponieważ pozytywnie wpływa on na odbiór estetyczny autobusu i kreuje przyjazną atmosferę dla podróżnych, szczególnie tych najmłodszych.
\MyFigure[0.4\textwidth]{images/solaris_jamnik.jpg}{Maskotka firmy Solaris}{fig:solarisjamnik}

Poniżej przedstawiono charakterystykę wszystkich wytypowanych modeli autobusów wraz z merytorycznym uzasadnieniem przyznanych im not.

\subsection{Solaris Urbino 15 LE electric}
Decydenci wysoko ocenili nowoczesny design zewnętrza pojazdu, oraz minimalistyczny projekt przestrzeni pasażerskiej. Znaczącym atutem pojazdu w kategorii komfortu są ergonomicznie wyprofilowane fotele, oferujące stabilne podparcie kręgosłupa podczas podróży.
\MyFigure[0.5\textwidth]{images/autobusy/Urbino15LEelectric/wnetrze.jpg}{Wnętrze autobusu Urbino 15 LE electric}{fig:Urbino15LEelectricwnetrze}
\MyFigure[0.5\textwidth]{images/autobusy/Urbino15LEelectric/zewnatrz.png}{Zewnętrze autobusu Urbino 15 LE electric}{fig:Urbino15LEelectriczewnatrz}

\subsection{CapaCity L}
Wysoka nota wynika przede wszystkim z doskonale zaprojektowanego wnętrza. Decydenci zwrócili szczególną uwagę na przestronność oraz imponującą liczbę dostępnych foteli. Całokształt aranżacji sprawia, że standard przejazdu przypomina podróż komfortowym autokarem wycieczkowym, co bezpośrednio przełoży się na wysoką satysfakcję pasażerów.
\MyFigure[0.5\textwidth]{images/autobusy/CapaCityL/wnetrze.png}{Wnętrze autobusu CapaCity L}{fig:CapaCityLwnetrze}
\MyFigure[0.5\textwidth]{images/autobusy/CapaCityL/zewnatrz.png}{Zewnętrze autobusu CapaCity L}{fig:CapaCityLzewnatrz}

\subsection{Citaro K}
Wysoka nota przyznana temu pojazdowi wynika w dużej mierze z bardzo udanego designu nadwozia. Front autobusu odznacza się masywną, solidną sylwetką, która budzi poczucie trwałości i bezpieczeństwa. Linia boczna harmonijnie współgra z estetyczną kolorystyką powłoki lakierniczej (w wariancie poglądowym jest to elegancki, złoty odcień). Kolejnym kluczowym atutem jest przemyślana aranżacja wnętrza. Zespół decyzyjny docenił układ tylnych rzędów siedzeń, których oparcia posiadają lekkie pochylenie, gwarantując pasażerom wyższy poziom odprężenia i komfortu w trakcie podróży.
\MyFigure[0.5\textwidth]{images/autobusy/CitaroK/wnetrze.png}{Wnętrze autobusu Citaro K}{fig:CitaroKwnetrze}
\MyFigure[0.5\textwidth]{images/autobusy/CitaroK/zewnatrz.png}{Zewnętrze autobusu Citaro K}{fig:CitaroKzewnatrz}

\subsection{Crossway 10.8 LE}
Niezwykle wysoka nota przyznana temu pojazdowi jest wynikiem jego wyjątkowych walorów estetycznych i funkcjonalnych. Zespół decyzyjny docenił przede wszystkim nowoczesną, modernistyczną linię nadwozia. Przestrzeń pasażerska wyróżnia się wysoką spójnością kolorystyczną oraz ogólną elegancją aranżacji. Na szczególne uznanie zasługują praktyczne udogodnienia w postaci dedykowanych przestrzeni bagażowych nad siedzeniami. Ponadto, zastosowanie podwyższonego podestu w tylnej sekcji pasażerskiej nadaje tej strefie ekskluzywny charakter, co znacząco podnosi prestiż i odczuwalny komfort podróży.
\MyFigure[0.5\textwidth]{images/autobusy/Crossway10.8LE/wnetrze.png}{Wnętrze autobusu Crossway 10.8 LE}{fig:Crossway108LEwnetrze}
\MyFigure[0.5\textwidth]{images/autobusy/Crossway10.8LE/zewnatrz.png}{Zewnętrze autobusu Crossway 10.8 LE}{fig:Crossway108LEzewnatrz}

\subsection{Crossway Pro 10,8}
Pojazd ten został oceniony na 6 punktów ze względu na swój surowy charakter. Decydenci określili design zewnętrzny jako nijaki i przestarzały, co w zestawieniu z proporcjami okien nadaje autobusowi ciężki, wizualnie zamknięty wygląd. Wnętrze pojazdu kontynuuje ten sam trend - aranżacja siedzeń, ułożonych w ścisłych, wysokich rzędach, sprawia wrażenie przytłaczającej i brakuje w niej elementów ocieplających wizerunek. W efekcie podróżni mogą odczuwać brak przestrzeni i swobody.
\MyFigure[0.5\textwidth]{images/autobusy/CrosswayPro10,8/wnetrze.png}{Wnętrze autobusu Crossway Pro 10,8}{fig:CrosswayPro108wnetrze}
\MyFigure[0.5\textwidth]{images/autobusy/CrosswayPro10,8/zewnatrz.png}{Zewnętrze autobusu Crossway Pro 10,8}{fig:CrosswayPro108zewnatrz}

\subsection{MAN Lion's City 10 E}
Pojazd ten otrzymał niższą punktację głównie ze względu na dysonans między poprawnym designem nadwozia a dyskusyjnym projektem przestrzeni pasażerskiej. Decydenci zwrócili szczególną uwagę na zbyt wąskie siedziska i oparcia foteli, które nie gwarantują odpowiedniego komfortu i mogą powodować uciążliwy ścisk w ułożeniu dwurzędowym. Dodatkowym czynnikiem negatywnie rzutującym na estetykę wnętrza jest nietypowe wykończenie podłogi. Wzór przypominający tradycyjne, domowe panele podłogowe został oceniony jako rozwiązanie nienaturalne i zaburzające profesjonalny wizerunek transportu zbiorowego.
\MyFigure[0.5\textwidth]{images/autobusy/MANLion'sCity10E/wnetrze.png}{Wnętrze autobusu MAN Lion's City 10 E}{fig:MANLionsCity10Ewnetrze}
\MyFigure[0.5\textwidth]{images/autobusy/MANLion'sCity10E/zewnatrz.png}{Zewnętrze autobusu MAN Lion's City 10 E}{fig:MANLionsCity10Ezewnatrz}

\subsection{MAN Lion's City 12 E}
Model ten zachowuje atrakcyjnie zaprojektowaną bryłę zewnętrzną, docenioną już w pokrewnych pojazdach tego producenta. Jego zdecydowaną przewagą jest natomiast znacznie przestronniejsze i bardziej funkcjonalne wnętrze. Zespół decyzyjny zauważył, że dodatkowa przestrzeń w kabinie pasażerskiej skutecznie poprawia ergonomię, oferując podróżnym o wiele wyższy poziom swobody oraz lepsze wrażenia z jazdy.
\MyFigure[0.5\textwidth]{images/autobusy/MANLion'sCity12E/wnetrze.png}{Wnętrze autobusu MAN Lion's City 12 E}{fig:MANLionsCity12Ewnetrze}
\MyFigure[0.5\textwidth]{images/autobusy/MANLion'sCity12E/zewnatrz.png}{Zewnętrze autobusu MAN Lion's City 12 E}{fig:MANLionsCity12Ezewnatrz}

\subsection{MAN Lion's City 12 G}
Niska ocena prestiżu jest w tym przypadku konsekwencją przeciętnej estetyki nadwozia oraz istotnych braków w wyposażeniu. Design zewnętrzny, choć w pełni akceptowalny, opiera się na powszechnie stosowanych schematach. Głównym powodem krytycznej oceny decydentów jest jednak wnętrze pojazdu. Kabina pasażerska charakteryzuje się niekorzystnym, surowym układem przestrzennym. Zainstalowane siedzenia wyglądają na niedostatecznie wyprofilowane, co negatywnie rzutuje na postrzegany komfort podróży. Dodatkowo, dużą wadą w tej kategorii okazał się brak zintegrowanych paneli wizualnej informacji dla podróżnych, które stanowią obecnie bazowy standard rynkowy.
\MyFigure[0.5\textwidth]{images/autobusy/MANLion'sCity12G/wnetrze.png}{Wnętrze autobusu MAN Lion's City 12 G}{fig:MANLionsCity12Gwnetrze}
\MyFigure[0.5\textwidth]{images/autobusy/MANLion'sCity12G/zewnatrz.png}{Zewnętrze autobusu MAN Lion's City 12 G}{fig:MANLionsCity12Gzewnatrz}

\subsection{MAN Lion's Intercity LE 14}
Model ten uzyskał umiarkowaną ocenę na poziomie 6 punktów. Choć ogólna linia nadwozia jest w dużej mierze poprawna, zespół decyzyjny krytycznie ocenił charakterystyczny, wydłużony pas przedni pojazdu, który wizualnie zaburza proporcje bryły zewnętrznej. Wnętrze autobusu wzbudziło mieszane odczucia. Z jednej strony pozytywnie oceniono podwyższoną sekcję tylną, urozmaicającą przestrzeń pasażerską. Z drugiej jednak, ocenę obniżyło ponowne zastosowanie dyskusyjnego wykończenia podłogi imitującego domowe panele drewniane. Istotnym zarzutem z perspektywy ergonomii jest również zbyt gęsty rozstaw foteli. Zredukowana odległość między rzędami znacząco ogranicza przestrzeń na nogi, co może powodować dyskomfort, w szczególności u osób o wyższym wzroście.
\MyFigure[0.4\textwidth]{images/autobusy/MANLion'sIntercityLE14/wnetrze.png}{Wnętrze autobusu MAN Lion's Intercity LE 14}{fig:MANLionsIntercityLE14wnetrze}
\MyFigure[0.5\textwidth]{images/autobusy/MANLion'sIntercityLE14/zewnatrz.png}{Zewnętrze autobusu MAN Lion's Intercity LE 14}{fig:MANLionsIntercityLE14zewnatrz}

\subsection{Nesobus 12}
Prezentowany model to niekwestionowany faworyt zespołu decyzyjnego w kategorii prestiżu. Zewnętrzna bryła pojazdu przyciąga uwagę modernistycznym, niezwykle nowoczesnym designem. Kluczowym czynnikiem, który zadecydował o tak wysokiej nocie, jest jednak bezprecedensowa jakość aranżacji wnętrza. Przestrzeń pasażerska odznacza się elegancją oraz harmonijnie dobraną paletą barw. Co istotne, zastosowana tu wykładzina drewnopodobna – w przeciwieństwie do innych ocenianych modeli – nie sprawia wrażenia taniego wykończenia. Wraz z architekturą siedzeń kreuje ona wysoce prestiżową, wręcz luksusową atmosferę. Na szczególne uznanie zasługuje innowacyjne zagospodarowanie strefy dla pasażerów stojących, obejmujące rozkładane fotele oraz tapicerowane, miękkie oparcia. Całości dopełnia wysoce ergonomiczny i optymalny układ standardowych miejsc siedzących.
\MyFigure[0.5\textwidth]{images/autobusy/Nesobus12/wnetrze.png}{Wnętrze autobusu Nesobus 12}{fig:Nesobus12wnetrze}
\MyFigure[0.5\textwidth]{images/autobusy/Nesobus12/zewnatrz.jpg}{Zewnętrze autobusu Nesobus 12}{fig:Nesobus12zewnatrz}

\subsection{SANCITY 9LE}
Model ten uzyskał od zespołu decyzyjnego stosunkowo niską ocenę na poziomie 4 punktów. O ile stylistyka zewnętrzna pojazdu prezentuje się estetycznie i nie budzi większych zastrzeżeń, o tyle wnętrze spotkało się z surową krytyką. Przestrzeń pasażerska sprawia wrażenie monotonnej, wizualnie pustej i pozbawionej wyrazu. Zauważalna jest nieoptymalna aranżacja kabiny, która skutkuje wyraźnym i nieuzasadnionym deficytem miejsc siedzących względem dostępnej kubatury. Jedynym elementem architektonicznym, który w pewnym stopniu urozmaica to surowe wnętrze i zapobiega przyznaniu jeszcze niższej noty, jest zastosowanie podwyższenia w tylnej sekcji pojazdu.
\MyFigure[0.5\textwidth]{images/autobusy/SANCITY9LE/wnetrze.png}{Wnętrze autobusu SANCITY 9LE}{fig:SANCITY9LEwnetrze}
\MyFigure[0.5\textwidth]{images/autobusy/SANCITY9LE/zewnatrz.png}{Zewnętrze autobusu SANCITY 9LE}{fig:SANCITY9LEzewnatrz}

\subsection{SANCITY 10LF}
Miejsce na opis i uzasadnienie oceny prestiżu dla tego modelu...
\MyFigure[0.5\textwidth]{images/autobusy/SANCITY10LF/wnetrze.png}{Wnętrze autobusu SANCITY 10LF}{fig:SANCITY10LFwnetrze}
\MyFigure[0.5\textwidth]{images/autobusy/SANCITY10LF/zewnatrz.png}{Zewnętrze autobusu SANCITY 10LF}{fig:SANCITY10LFzewnatrz}

\subsection{SANCITY 10LFE}
Miejsce na opis i uzasadnienie oceny prestiżu dla tego modelu...
\MyFigure[0.4\textwidth]{images/autobusy/SANCITY10LFE/wnetrze.png}{Wnętrze autobusu SANCITY 10LFE}{fig:SANCITY10LFEwnetrze}
\MyFigure[0.5\textwidth]{images/autobusy/SANCITY10LFE/zewnatrz.png}{Zewnętrze autobusu SANCITY 10LFE}{fig:SANCITY10LFEzewnatrz}

\subsection{SANCITY 12LF}
Miejsce na opis i uzasadnienie oceny prestiżu dla tego modelu...
\MyFigure[0.5\textwidth]{images/autobusy/SANCITY12LF/wnetrze.png}{Wnętrze autobusu SANCITY 12LF}{fig:SANCITY12LFwnetrze}
\MyFigure[0.5\textwidth]{images/autobusy/SANCITY12LF/zewnatrz.png}{Zewnętrze autobusu SANCITY 12LF}{fig:SANCITY12LFzewnatrz}

\subsection{SANCITY 12LFH}
Miejsce na opis i uzasadnienie oceny prestiżu dla tego modelu...
\MyFigure[0.5\textwidth]{images/autobusy/SANCITY12LFH/wnetrze.png}{Wnętrze autobusu SANCITY 12LFH}{fig:SANCITY12LFHwnetrze}
\MyFigure[0.5\textwidth]{images/autobusy/SANCITY12LFH/zewnatrz.png}{Zewnętrze autobusu SANCITY 12LFH}{fig:SANCITY12LFHzewnatrz}

\subsection{Urbino 8,9 LE}
Miejsce na opis i uzasadnienie oceny prestiżu dla tego modelu...
\MyFigure[0.5\textwidth]{images/autobusy/Urbino8,9LE/wnetrze.png}{Wnętrze autobusu Urbino 8,9 LE}{fig:Urbino89LEwnetrze}
\MyFigure[0.5\textwidth]{images/autobusy/Urbino8,9LE/zewnatrz.png}{Zewnętrze autobusu Urbino 8,9 LE}{fig:Urbino89LEzewnatrz}

\subsection{Urbino 8,9 LE electric}
Miejsce na opis i uzasadnienie oceny prestiżu dla tego modelu...
\MyFigure[0.5\textwidth]{images/autobusy/Urbino8,9LEelectric/wnetrze.png}{Wnętrze autobusu Urbino 8,9 LE electric}{fig:Urbino89LEelectricwnetrze}
\MyFigure[0.5\textwidth]{images/autobusy/Urbino8,9LEelectric/zewnatrz.png}{Zewnętrze autobusu Urbino 8,9 LE electric}{fig:Urbino89LEelectriczewnatrz}

\subsection{Urbino 10,5 electric}
Miejsce na opis i uzasadnienie oceny prestiżu dla tego modelu...
\MyFigure[0.5\textwidth]{images/autobusy/Urbino10,5electric/wnetrze.png}{Wnętrze autobusu Urbino 10,5 electric}{fig:Urbino105electricwnetrze}
\MyFigure[0.5\textwidth]{images/autobusy/Urbino10,5electric/zewnatrz.jpg}{Zewnętrze autobusu Urbino 10,5 electric}{fig:Urbino105electriczewnatrz}

\subsection{Urbino 12}
Miejsce na opis i uzasadnienie oceny prestiżu dla tego modelu...
\MyFigure[0.5\textwidth]{images/autobusy/Urbino12/wnetrze.png}{Wnętrze autobusu Urbino 12}{fig:Urbino12wnetrze}
\MyFigure[0.5\textwidth]{images/autobusy/Urbino12/zewnatrz.jpg}{Zewnętrze autobusu Urbino 12}{fig:Urbino12zewnatrz}

\subsection{Urbino 12 CNG}
Miejsce na opis i uzasadnienie oceny prestiżu dla tego modelu...
\MyFigure[0.5\textwidth]{images/autobusy/Urbino12CNG/wnetrze.png}{Wnętrze autobusu Urbino 12 CNG}{fig:Urbino12CNGwnetrze}
\MyFigure[0.5\textwidth]{images/autobusy/Urbino12CNG/zewnatrz.png}{Zewnętrze autobusu Urbino 12 CNG}{fig:Urbino12CNGzewnatrz}

\subsection{Urbino 15 LE electric}
Miejsce na opis i uzasadnienie oceny prestiżu dla tego modelu...
\MyFigure[0.5\textwidth]{images/autobusy/Urbino15LEelectric/wnetrze.jpg}{Wnętrze autobusu Urbino 15 LE electric}{fig:Urbino15LEelectricwnetrze}
\MyFigure[0.5\textwidth]{images/autobusy/Urbino15LEelectric/zewnatrz.png}{Zewnętrze autobusu Urbino 15 LE electric}{fig:Urbino15LEelectriczewnatrz}

\subsection{Urbino 18}
Miejsce na opis i uzasadnienie oceny prestiżu dla tego modelu...
\MyFigure[0.5\textwidth]{images/autobusy/Urbino18/wnetrze.png}{Wnętrze autobusu Urbino 18}{fig:Urbino18wnetrze}
\MyFigure[0.5\textwidth]{images/autobusy/Urbino18/zewnatrz.jpg}{Zewnętrze autobusu Urbino 18}{fig:Urbino18zewnatrz}

\subsection{Urbino 18 CNG}
Miejsce na opis i uzasadnienie oceny prestiżu dla tego modelu...
\MyFigure[0.5\textwidth]{images/autobusy/Urbino18CNG/wnetrze.png}{Wnętrze autobusu Urbino 18 CNG}{fig:Urbino18CNGwnetrze}
\MyFigure[0.7\textwidth]{images/autobusy/Urbino18CNG/zewnatrz.png}{Zewnętrze autobusu Urbino 18 CNG}{fig:Urbino18CNGzewnatrz}

\section{Analiza SAW dla autobusów dużych}

Do tej kategorii zakwalifikowano pojazdy o wysokiej zdolności przewozowej, przeznaczone na główne osie komunikacyjne miasta.

\begin{table}[H]
    \centering
    \caption{Macierz decyzyjna (dane wejściowe) -- autobusy duże}
    \vspace{0.2cm}
    \begin{adjustbox}{max width=\textwidth}
    \begin{tabular}{|l|c|c|c|c|c|}
        \hline
        \textbf{Model Autobusu} & \textbf{K1: Moc [KM]} & \textbf{K2: Miejsca} & \textbf{K3: Prestiż} & \textbf{K4: Udogod.} & \textbf{K5: Koszt [mln]} \\ 
        & \textbf{(MAX)} & \textbf{(MAX)} & \textbf{(MAX)} & \textbf{(MAX)} & \textbf{(MIN)} \\ \hline
        Urbino 15 LE electric & 408 & 65 & 7.00 & 2 & 2.20 \\ \hline
        Urbino 18 CNG & 320 & 55 & 10.00 & 3 & 2.40 \\ \hline
        Urbino 18 & 310 & 54 & 6.33 & 3 & 1.80 \\ \hline
        Urbino 12 & 256 & 16 & 8.67 & 3 & 1.00 \\ \hline
        Urbino 12 CNG & 320 & 40 & 4.67 & 3 & 1.10 \\ \hline
        MAN Lion's City 12 E & 326 & 44 & 6.33 & 4 & 2.60 \\ \hline
        MAN Lion's Intercity LE 14 & 265 & 63 & 6.33 & 3 & 1.80 \\ \hline
        MAN Lion's City 12 G & 280 & 37 & 3.00 & 4 & 2.00 \\ \hline
        SANCITY 12LFH & 344 & 40 & 8.00 & 3 & 3.25 \\ \hline
        SANCITY 12LF & 277 & 40 & 4.00 & 3 & 1.00 \\ \hline
        CapaCity L & 360 & 44 & 8.67 & 4 & 2.00 \\ \hline
        Nesobus 12 & 340 & 37 & 9.67 & 3 & 2.90 \\ \hline
    \end{tabular}
    \end{adjustbox}
\end{table}

\begin{table}[H]
    \centering
    \caption{Wyniki modelu SAW oraz ranking -- autobusy duże}
    \vspace{0.2cm}
    \begin{adjustbox}{max width=\textwidth}
    \begin{tabular}{|c|l|c|c|c|c|c|c|}
        \hline
        \textbf{Poz.} & \textbf{Model Autobusu} & \textbf{K1 norm.} & \textbf{K2 norm.} & \textbf{K3 norm.} & \textbf{K4 norm.} & \textbf{K5 norm.} & \textbf{Wynik SAW} \\ \hline
        1 & Urbino 18 CNG & 0.78 & 0.85 & 1.00 & 0.75 & 0.42 & \textbf{0.812} \\ \hline
        2 & Urbino 15 LE electric & 1.00 & 1.00 & 0.70 & 0.50 & 0.45 & \textbf{0.807} \\ \hline
        3 & CapaCity L & 0.88 & 0.68 & 0.87 & 1.00 & 0.50 & \textbf{0.796} \\ \hline
        4 & Nesobus 12 & 0.83 & 0.57 & 0.97 & 0.75 & 0.34 & \textbf{0.736} \\ \hline
        5 & MAN Lion's Intercity LE 14 & 0.65 & 0.97 & 0.63 & 0.75 & 0.56 & \textbf{0.729} \\ \hline
        6 & Urbino 18 & 0.76 & 0.83 & 0.63 & 0.75 & 0.56 & \textbf{0.722} \\ \hline
        7 & SANCITY 12LFH & 0.84 & 0.62 & 0.80 & 0.75 & 0.31 & \textbf{0.704} \\ \hline
        8 & MAN Lion's City 12 E & 0.80 & 0.68 & 0.63 & 1.00 & 0.38 & \textbf{0.701} \\ \hline
        9 & Urbino 12 CNG & 0.78 & 0.62 & 0.47 & 0.75 & 0.91 & \textbf{0.668} \\ \hline
        10 & Urbino 12 & 0.63 & 0.25 & 0.87 & 0.75 & 1.00 & \textbf{0.645} \\ \hline
        11 & SANCITY 12LF & 0.68 & 0.62 & 0.40 & 0.75 & 1.00 & \textbf{0.633} \\ \hline
        12 & MAN Lion's City 12 G & 0.69 & 0.57 & 0.30 & 1.00 & 0.50 & \textbf{0.569} \\ \hline
    \end{tabular}
    \end{adjustbox}
\end{table}

Skupienie wagi na prestiżu i ergonomii pasażerskiej doprowadziło do zwycięstwa pojazdu z napędem gazowym -- \textit{Solaris Urbino 18 CNG}. Zgromadził on doskonałe noty za estetykę i oferowane rozwiązania (prestiż uśredniony na 10,00) oraz dużą pojemność. Bezpośrednio za nim znalazł się najpotężniejszy (408 KM) elektryczny \textit{Urbino 15 LE}.

\section{Analiza SAW dla autobusów średnich}

Zestawienie pojazdów o długości do kilkunastu metrów, które idealnie wpisują się w linie dowozowe, osiedlowe oraz trasy o węższej infrastrukturze drogowej.

\begin{table}[H]
    \centering
    \caption{Macierz decyzyjna (dane wejściowe) -- autobusy średnie}
    \vspace{0.2cm}
    \begin{adjustbox}{max width=\textwidth}
    \begin{tabular}{|l|c|c|c|c|c|}
        \hline
        \textbf{Model Autobusu} & \textbf{K1: Moc [KM]} & \textbf{K2: Miejsca} & \textbf{K3: Prestiż} & \textbf{K4: Udogod.} & \textbf{K5: Koszt [mln]} \\ 
        & \textbf{(MAX)} & \textbf{(MAX)} & \textbf{(MAX)} & \textbf{(MAX)} & \textbf{(MIN)} \\ \hline
        Urbino 10,5 electric & 340 & 33 & 4.00 & 3 & 2.60 \\ \hline
        Urbino 8,9 LE electric & 218 & 27 & 8.00 & 4 & 2.40 \\ \hline
        Urbino 8,9 LE & 250 & 33 & 6.67 & 4 & 1.10 \\ \hline
        MAN Lion's City 10 E & 326 & 36 & 5.00 & 4 & 2.60 \\ \hline
        SANCITY 9LE & 210 & 16 & 4.33 & 4 & 0.60 \\ \hline
        SANCITY 10LFE & 340 & 24 & 5.00 & 4 & 2.20 \\ \hline
        SANCITY 10LF & 247 & 26 & 5.67 & 3 & 0.90 \\ \hline
        Citaro K & 299 & 32 & 8.33 & 5 & 1.30 \\ \hline
        Crossway 10.8LE & 330 & 29 & 8.67 & 4 & 1.10 \\ \hline
        Crossway Pro 10,8 & 360 & 47 & 6.33 & 3 & 0.90 \\ \hline
    \end{tabular}
    \end{adjustbox}
\end{table}

\begin{table}[H]
    \centering
    \caption{Wyniki modelu SAW oraz ranking -- autobusy średnie}
    \vspace{0.2cm}
    \begin{adjustbox}{max width=\textwidth}
    \begin{tabular}{|c|l|c|c|c|c|c|c|}
        \hline
        \textbf{Poz.} & \textbf{Model Autobusu} & \textbf{K1 norm.} & \textbf{K2 norm.} & \textbf{K3 norm.} & \textbf{K4 norm.} & \textbf{K5 norm.} & \textbf{Wynik SAW} \\ \hline
        1 & Crossway Pro 10,8 & 1.00 & 1.00 & 0.73 & 0.60 & 0.67 & \textbf{0.849} \\ \hline
        2 & Crossway 10.8LE & 0.92 & 0.62 & 1.00 & 0.80 & 0.55 & \textbf{0.807} \\ \hline
        3 & Citaro K & 0.83 & 0.68 & 0.96 & 1.00 & 0.46 & \textbf{0.804} \\ \hline
        4 & Urbino 8,9 LE & 0.69 & 0.70 & 0.77 & 0.80 & 0.55 & \textbf{0.711} \\ \hline
        5 & MAN Lion's City 10 E & 0.91 & 0.77 & 0.58 & 0.80 & 0.23 & \textbf{0.698} \\ \hline
        6 & Urbino 8,9 LE electric & 0.61 & 0.57 & 0.92 & 0.80 & 0.25 & \textbf{0.662} \\ \hline
        7 & SANCITY 10LFE & 0.94 & 0.51 & 0.58 & 0.80 & 0.27 & \textbf{0.646} \\ \hline
        8 & Urbino 10,5 electric & 0.94 & 0.70 & 0.46 & 0.60 & 0.23 & \textbf{0.639} \\ \hline
        9 & SANCITY 10LF & 0.69 & 0.55 & 0.65 & 0.60 & 0.67 & \textbf{0.632} \\ \hline
        10 & SANCITY 9LE & 0.58 & 0.34 & 0.50 & 0.80 & 1.00 & \textbf{0.568} \\ \hline
    \end{tabular}
    \end{adjustbox}
\end{table}

W tym segmencie dominację pokazała marka Iveco za sprawą modelu \textit{Crossway Pro 10,8}. Oferuje on absolutnie bezkonkurencyjną moc silnika w swojej klasie gabarytowej (360 KM) połączoną z maksymalizacją wykorzystania wnętrza pojazdu (najwyższa ocena w parametrze "Miejsca"). Świetnie wyważony pod kątem prestiżu w stosunku do mocy okazał się też \textit{Mercedes-Benz Citaro K}, zyskując najwyższą notę za udogodnienia dla niepełnosprawnych.